\FloatBarrier
\section{Limitations}
\label{sec:limitations}

\subsection{Where \colbin{} loses or ties}

\paragraph{Single records, against Avro.} On records with no table in
them \colbin{} is a row format with one byte of header per field, and Avro,
which writes no header at all, is smaller as written: on SchemaStore
\colbin{} against Avro is \result{vs.datasets.schemastore.avro.raw.factor}
as written and \result{vs.datasets.schemastore.avro.gzip.factor} under
gzip, smaller in
\result{vs.datasets.schemastore.avro.gzip.wins} of
\result{vs.datasets.schemastore.avro.gzip.n} documents. Avro pays for this
with a reader that must have the writer's exact schema.

\paragraph{Time on single records, against generated code.} On SchemaStore
\colbin{}'s encode time against \fmt{vtprotobuf}'s is
\result{vs.datasets.schemastore.vtprotobuf.encode.factor} and against
generated MessagePack's
\result{vs.datasets.schemastore.msgpack.encode.factor}; it encodes
faster in \result{vs.datasets.schemastore.vtprotobuf.encode.wins} of
\result{vs.datasets.schemastore.vtprotobuf.encode.n} and
\result{vs.datasets.schemastore.msgpack.encode.wins} of
\result{vs.datasets.schemastore.msgpack.encode.n} documents, and to decode it is
\result{vs.datasets.schemastore.msgpack.decode-fresh.factor} against
generated MessagePack and
\result{vs.datasets.schemastore.flatbuffers.decode-fresh.factor} against
FlatBuffers. Straight-line generated code beats a plan interpreted
per field on a message of a few fields; the columnar path is what makes up
the difference on tables.

\paragraph{Access without decoding.} \colbin{} must decode a message to read
one field of it. Arrow and FlatBuffers read a column in place and are faster
at it by orders of magnitude (Section~\ref{sec:results}).

\paragraph{Data with nothing to exploit.} A column of random 64-bit values
costs its eight bytes a value plus a width byte per block; strings are not
compressed at all by the default format. A table that is mostly long,
distinct strings is the size of its strings in any format.

\subsection{What the evaluation does not cover}

\paragraph{The author evaluates his own format.} The format, the harness and
the choices below are one person's, and a reader should weigh the results
accordingly. Four things bound the choices. The data is third-party, and the
Go types of \result{cases.categories.bulk} bulk and
\result{cases.categories.web-tables} web tables are generated from the
sources' \code{CREATE TABLE} statements; only SchemaStore's are written by
hand (below). Every other format's schema is generated from those same
types, so no format is given a type of its own. Each library is called with
its default options and through the interface its documentation gives for a
hot path: generated code encodes into the caller's buffer, the FlatBuffers
builder is reused, \colbin{} runs through a codec handle built once per
type, and converting between a generated type and the Go struct is not
timed. And the departures from a default are few and stated here: CBOR's
decoder has its limits on array and map sizes raised, which a table of
$10^5$ rows exceeds; Arrow's buffers are not compressed, so that its size as
written is not a compressed one; \fmt{gob} writes one message per encoder,
and so describes its types in every message. Two choices remain that favour
\colbin{}. The rule for types, decimals and money as integer minor units and
dates as integer days, keeps the sources' exact types exact, but it puts
most values where an integer codec works, which is \colbin{}'s strength; the
evaluation does not measure a schema that keeps them as floats or strings.
And the
libraries, one or two per format, were chosen by the author as the ones in
common use in Go; another implementation of the same format would time
differently.

\paragraph{One language and one implementation per format.} All measurements
are of Go implementations. A format's speed is its implementation's: the
two implementations of Protocol Buffers and of MessagePack measured here
differ by more than most pairs of formats do. The sizes carry over to other
languages; the times do not.

\paragraph{Two machines, both x86-64 and of one vendor.} The main times come from one server
processor model, in four cloud container instances of the same shape (one
per dataset), whose clock the host controls and whose cache and memory
bandwidth other tenants share. A summary over a regime pools cases measured
in different instances; a ratio on one case never does. We measured that
clock rather than fixed it (Section~\ref{sec:machine}). Absolute times there
were not repeatable to 3\,\% between two passes; the ratios this paper
reports were, to within \result{repeat.ratio-worst}\,\%, and differences
below about 4\,\% should not be read as differences. The second machine, a
laptop with its clock fixed (Section~\ref{sec:second-machine}), shows which
conclusions survive a change of processor, but it is a single unit of a
single model, and two machines are not a population. Both are x86-64
processors of AMD, one generation apart: no Intel processor and no arm64
machine was measured.

\paragraph{Datasets not covered.} Real floats and timestamps are
represented by Public BI only. No dataset here was assembled by a third
party specifically to compare serialization formats except SchemaStore's
corpus, and the evaluation is not validated against an existing third-party
benchmark suite.

\paragraph{Synthetic distributions.} TPC-H and TPC-DS draw values from
near-uniform distributions, which under-represent the skew and correlation
of real data~\cite{vogelsgesang2018getreal}. They are reported per dataset
so that Public BI, which is real, can be read on its own
(Section~\ref{sec:by-dataset}).

\paragraph{Hand-written types for SchemaStore.} The Go structs for the 27
SchemaStore documents were written by the author, from the corpus's
Protocol Buffers schemas. Choices of integer width, of pointer or value, and
of field numbering affect every schema-driven format, not only \colbin{},
but they are choices. The artifact also measures each document as a
\code{map[string]any}, which involves none; only the schemaless formats and
Protocol Buffers (through \code{google.protobuf.Value}) can carry that, so it
compares \colbin{} with a few formats on cases the others do not have, and
it is not reported here.
\result{coverage.web.colbin.na} documents cannot be given a type \colbin{}
carries and are reported as not applicable, which removes them from
\colbin{}'s comparisons; the other formats' means against JSON
(Appendix~\ref{sec:against-json}) are over them.

\paragraph{gob.} \fmt{gob} omits zero values and does not distinguish a nil
pointer from a pointer to zero, so it does not round-trip a table with a
nullable column that holds a zero:
\result{coverage.web.gob.lossy} of \result{cases.web} web cases and
\result{coverage.bulk.gob.lossy} of \result{cases.bulk} bulk cases. On those
it is sized but not timed, so its timing means are over fewer cases, and
they are the cases without nullable columns. Its sizes also depend on what
the process encoded before, because \fmt{gob} assigns type identifiers per
process in order of first use; all cases of a run are sized in one process,
in a fixed order, and a different selection of cases gives sizes that
differ by up to a dozen bytes a message. Both machines sized one dataset
per process, so their sizes agree on every case reported here.

\paragraph{Stability of the format.} \colbin{} is at version
\result{machine.colbin-version}, before 1.0, with no compatibility promise
between minor versions. The numbers are for that version.

\paragraph{Not measured.} The cost of the first use of a type, when the
encoding plan is built, is not measured: every time in this paper is of a
warmed-up process. Memory is reported for decode only
(Section~\ref{sec:memory}), as what an operation allocates, not as what the
process holds. Schema evolution is described and tested for behaviour
(Section~\ref{sec:evolution}), not for cost.
